%2multibyte Version: 5.50.0.2953 CodePage: 1252

\documentclass{article}
%%%%%%%%%%%%%%%%%%%%%%%%%%%%%%%%%%%%%%%%%%%%%%%%%%%%%%%%%%%%%%%%%%%%%%%%%%%%%%%%%%%%%%%%%%%%%%%%%%%%%%%%%%%%%%%%%%%%%%%%%%%%%%%%%%%%%%%%%%%%%%%%%%%%%%%%%%%%%%%%%%%%%%%%%%%%%%%%%%%%%%%%%%%%%%%%%%%%%%%%%%%%%%%%%%%%%%%%%%%%%%%%%%%%%%%%%%%%%%%%%%%%%%%%%%%%
\usepackage{amsfonts}
\usepackage{amsmath}

\setcounter{MaxMatrixCols}{10}
%TCIDATA{OutputFilter=LATEX.DLL}
%TCIDATA{Version=5.50.0.2953}
%TCIDATA{Codepage=1252}
%TCIDATA{<META NAME="SaveForMode" CONTENT="1">}
%TCIDATA{BibliographyScheme=Manual}
%TCIDATA{Created=Tuesday, September 13, 2016 11:26:35}
%TCIDATA{LastRevised=Sunday, November 03, 2019 21:33:40}
%TCIDATA{<META NAME="GraphicsSave" CONTENT="32">}
%TCIDATA{<META NAME="DocumentShell" CONTENT="Standard LaTeX\Blank - Standard LaTeX Article">}
%TCIDATA{Language=American English}
%TCIDATA{CSTFile=40 LaTeX article.cst}

\newtheorem{theorem}{Theorem}
\newtheorem{acknowledgement}[theorem]{Acknowledgement}
\newtheorem{algorithm}[theorem]{Algorithm}
\newtheorem{axiom}[theorem]{Axiom}
\newtheorem{case}[theorem]{Case}
\newtheorem{claim}[theorem]{Claim}
\newtheorem{conclusion}[theorem]{Conclusion}
\newtheorem{condition}[theorem]{Condition}
\newtheorem{conjecture}[theorem]{Conjecture}
\newtheorem{corollary}[theorem]{Corollary}
\newtheorem{criterion}[theorem]{Criterion}
\newtheorem{definition}[theorem]{Definition}
\newtheorem{example}[theorem]{Example}
\newtheorem{exercise}[theorem]{Exercise}
\newtheorem{lemma}[theorem]{Lemma}
\newtheorem{notation}[theorem]{Notation}
\newtheorem{problem}[theorem]{Problem}
\newtheorem{proposition}[theorem]{Proposition}
\newtheorem{remark}[theorem]{Remark}
\newtheorem{solution}[theorem]{Solution}
\newtheorem{summary}[theorem]{Summary}
\newenvironment{proof}[1][Proof]{\noindent\textbf{#1.} }{\ \rule{0.5em}{0.5em}}
\input{tcilatex}
\begin{document}


\begin{center}
{\Large Class Test}
\end{center}

Question 1

(a) Define the mean squared error (MSE) of an estimator $\hat{\theta}_{n}$.
Show how it relates to the bias and variance of the estimator, when $\dim
\left( \theta \right) =1$. [10 marks]

\begin{definition}
\textbf{The mean squared error} (MSE) of a statistic $\hat{\theta}_{n}$ is
defined as 
\begin{eqnarray*}
MSE\left( \hat{\theta}_{n},\theta \right) &:&=\mathbb{E}\left[ \left( \hat{%
\theta}_{n}-\theta \right) ^{2}\right] \\
&=&\mathbb{E}\left( \hat{\theta}_{n}-\mathbb{E}\hat{\theta}_{n}\right)
^{2}+\left( \mathbb{E}\hat{\theta}_{n}-\theta \right) ^{2}+2\left( \mathbb{E}%
[\hat{\theta}_{n}]-\theta \right) \mathbb{E}\left[ \left( \hat{\theta}_{n}-%
\mathbb{E}[\hat{\theta}_{n}]\right) \right] \\
&=&var\left( \hat{\theta}_{n}\right) +bias^{2}\left( \hat{\theta}_{n}\right)
\end{eqnarray*}
\end{definition}

(b) State Central Limit Theorem for a sequence of i.i.d. random variables
carefully stating any conditions required. Briefly explain the theorem's
importance for statistical inference. [10 marks]

\begin{theorem}[The Central Limit Theorem]
Let $X_{1},...,X_{n}$ be i.i.d. with $\mathbb{E}\left( X_{1}\right) =\mu $
and $V\left( X_{1}\right) =\sigma ^{2}<\infty .$ Then, 
\begin{eqnarray*}
Z_{n} &=&\frac{\overline{X}_{n}-\mu }{\sqrt{V(\overline{X}_{n})}}=\frac{%
\sqrt{n}}{\sigma }(\overline{X}_{n}-\mu )=\frac{1}{\sigma \sqrt{n}}%
\sum_{i=1}^{n}(X_{i}-\mu )\rightarrow _{d}Z\text{ } \\
\text{\ where }Z &\sim &\mathcal{N}(0,1)
\end{eqnarray*}

In other words, 
\begin{equation*}
\lim_{n\rightarrow \infty }F_{Z_{n}}=\Phi (z)=\int_{-\infty }^{z}(2\pi
)^{-1/2}\exp (-x^{2}/2)dx.
\end{equation*}
\end{theorem}

The CLT\ says that probability of $\frac{\sqrt{n}\left( \bar{X}_{n}-\mu
\right) }{\sigma }$ converges to a $N\left( 0,1\right) $\ r.v. This is a
remarkable result since nothing was assumed about the distribution of the
i.i.d. sequence $\left\{ X_{1},..,X_{n}\right\} $, apart from the existence
of their mean and variance; so the $X^{\prime }s$ could be generated by any
distribution. The theorem can be used for confidence interval construction
and hypothesis testing.

(c) Define a consistent estimator and explain why this is an important
property. [10 marks]

\begin{definition}
An estimator $\hat{\theta}_{n}$ of $\theta $ is called \textbf{consistent}
if $\hat{\theta}_{n}\rightarrow _{p}\theta $ as the sample size $n$ tends to 
$\infty $.
\end{definition}

Consistency is an important property which implies that the estimator gets
more precisely centered around the true value as more observations arrive.
At the limit, the distribution of $\hat{\theta}_{n}$ collapses to a point
and that point is the true unknown $\theta .$

Question 2

(a) Provide the definition for an efficient estimator. State and prove the
Gauss Markov theorem for $\hat{\beta}^{OLS}$. [15 marks]

\begin{definition}
An estimator $\hat{\theta}_{n}$ is called efficient if it is unbiased and
has the smallest variance, that is, $v[\widehat{\theta }_{n}]\leq v[\tilde{%
\theta}_{n}]$ for all other unbiased estimators $\tilde{\theta}_{n}.$
\end{definition}

\begin{theorem}[Gauss Markov]
Let $\widetilde{\beta }$ be any linear unbiased estimator for $\beta $ in
the linear regression model. Then, $var(\hat{\beta}^{OLS}|X)\leq var(%
\widetilde{\beta }|X).$
\end{theorem}

\begin{proof}
Consider another linear estimator $\widetilde{\beta }=A^{\prime }y$ where $A$
is an $n\times k$ function of $X.$ Since $\widetilde{\beta }$ is unbiased, 
\begin{equation*}
\beta =\mathbb{E[}\widetilde{\beta }|X]=A^{\prime }\mathbb{E}\left[ y|X%
\right] =A^{\prime }\mathbb{E}\left[ \left( X\beta +\varepsilon \right) |X%
\right] =A^{\prime }X\beta ,
\end{equation*}%
hence $A^{\prime }X=I_{k}.$ Moreover, 
\begin{equation*}
var\left( \widetilde{\beta }|X\right) =var\left( A^{\prime }y|X\right)
=A^{\prime }var\left( y|X\right) A=A^{\prime }var\left( \varepsilon
|X\right) A=\sigma ^{2}A^{\prime }A.
\end{equation*}%
We also know, that $var\left( \hat{\beta}|X\right) =\sigma ^{2}(X^{\prime
}X)^{-1}.$ Now, let's show that $var\left( \widetilde{\beta }|X\right) \geq
var\left( \hat{\beta}|X\right) ,$ that is $A^{\prime }A-(X^{\prime }X)^{-1}$
is positive semi-definite matrix. Recall that $A^{\prime }X=I_{k}$ and
define $C=A-X\left( X^{\prime }X\right) ^{-1},$ noting that $X^{\prime
}C=X^{\prime }A-X^{\prime }X\left( X^{\prime }X\right) ^{-1}=0.$ We have 
\begin{gather*}
A^{\prime }A-(X^{\prime }X)^{-1}=\left( C+X\left( X^{\prime }X\right)
^{-1}\right) ^{\prime }\left( C+X\left( X^{\prime }X\right) ^{-1}\right)
-(X^{\prime }X)^{-1} \\
=C^{\prime }C+C^{\prime }X\left( X^{\prime }X\right) ^{-1}+\left( X^{\prime
}X\right) ^{-1}X^{\prime }C+ \\
\left( X^{\prime }X\right) ^{-1}X^{\prime }X\left( X^{\prime }X\right)
^{-1}-(X^{\prime }X)^{-1}=C^{\prime }C\geq 0
\end{gather*}%
since $C^{\prime }C$ is a quadratic form, hence positive semi-definite. So, $%
\widehat{\beta }^{OLS}$ is efficient among a class of linear estimators.
\end{proof}

(b) Derive the asymptotic distribution of $\sqrt{n}\left( \hat{\beta}%
^{OLS}-\beta \right) $, carefully explaining each step of your derivations.
[15 marks]

\begin{proof}
$\sqrt{n}(\widehat{\beta }_{n}-\beta )=\left( \frac{1}{n}\tsum%
\nolimits_{i=1}^{n}\mathbf{x}_{i}\mathbf{x}_{i}^{\prime }\right) ^{-1}\frac{1%
}{\sqrt{n}}\tsum\nolimits_{i=1}^{n}\mathbf{x}_{i}\varepsilon _{i}$%
\begin{gather}
\left( \frac{1}{n}\tsum\nolimits_{i=1}^{n}\mathbf{x}_{i}\mathbf{x}%
_{i}^{\prime }\right) ^{-1}\rightarrow _{p}Q^{-1}\text{ {\tiny by WLLN, CMT
and Ass. 4}}  \notag \\
\frac{1}{\sqrt{n}}(\tsum\nolimits_{i=1}^{n}\mathbf{x}_{i}\varepsilon _{i}-%
\mathbb{E[}\mathbf{x}_{1}\varepsilon _{1}])\rightarrow _{d}\mathcal{N}(0,var(%
\mathbf{x}_{1}\varepsilon _{1}))\text{ {\tiny by the CLT}}  \notag \\
\frac{1}{\sqrt{n}}(\tsum\nolimits_{i=1}^{n}\mathbf{x}_{i}\varepsilon
_{i})\rightarrow _{d}\mathcal{N}(0,\sigma ^{2}Q)  \notag \\
\Rightarrow \sqrt{n}(\widehat{\beta }_{n}-\beta )\rightarrow _{d}\mathcal{N}%
(0,\sigma ^{2}Q^{-1}QQ^{-1})\text{ {\tiny by Slutsky Theorem}}  \notag \\
\sqrt{n}(\widehat{\beta }_{n}-\beta )\rightarrow _{d}\mathcal{N}(0,\sigma
^{2}Q^{-1}).  \label{clt}
\end{gather}%
since $\mathbb{E[}\mathbf{x}_{i}\varepsilon _{i}]=\mathbb{E[}\mathbf{x}%
_{1}\varepsilon _{1}]=0$ and $var(\mathbf{x}_{i}\varepsilon _{i})=var(%
\mathbf{x}_{1}\varepsilon _{1})=$ $\mathbb{E[}\mathbf{x}_{1}\mathbb{E}\left[
\varepsilon _{1}^{2}\right] \mathbf{x}_{1}^{\prime }]=\sigma ^{2}Q.$
\end{proof}

(c) Using your answer in (b), and supposing for simplicity that $\dim \left(
\beta \right) =1,$ derive the asymptotic distribution of $\sqrt{n}\left(
\exp \left( 7\hat{\beta}^{OLS}\right) -\exp \left( 7\beta \right) \right) $
and construct a 99\% confidence interval for $\exp \left( 7\beta \right) .$
[10 marks]

Applying the delta method, we immediately have that 
\begin{equation*}
\sqrt{n}\left( \exp \left( 7\hat{\beta}^{OLS}\right) -\exp \left( 7\beta
\right) \right) \rightarrow _{d}\mathcal{N}(0,49\exp \left( 14\beta \right)
\sigma ^{2}Q^{-1})
\end{equation*}

since $\frac{\partial \exp \left( 7\beta \right) }{\partial \beta }=7\exp
\left( 7\beta \right) .$ A 99\% confidence interval for $\exp \left( 7\beta
\right) $ is given by 
\begin{equation*}
\exp \left( 7\beta \right) \in \left( 7\hat{\beta}^{OLS}\pm 2.58\left( 7\exp
\left( 7\widehat{\beta }\right) \frac{1}{\sqrt{n}}\sqrt{\widehat{\sigma }^{2}%
\widehat{Q}^{-1}}\right) \right) \text{ with probability 99\%}.
\end{equation*}

Bonus points if student uses estimator for the variance and points out it is
consistent.

(d) Derive an expression for the bias of $\hat{\sigma}_{n}^{2}=\frac{1}{n}%
\hat{\varepsilon}^{\prime }\hat{\varepsilon}$. What happens to the bias as
the sample size $n$ increases? [15 marks] 
\begin{gather*}
\mathbb{E}[\widehat{\sigma }_{n}^{2}]=\mathbb{E}[\frac{1}{n}\widehat{%
\varepsilon }^{\prime }\widehat{\varepsilon }]=\mathbb{E}[\frac{1}{n}%
\varepsilon ^{\prime }M_{X}\varepsilon ]=\frac{1}{n}\mathbb{E}%
[tr(\varepsilon ^{\prime }M_{X}\varepsilon )]=\frac{1}{n}\mathbb{E}%
[tr(M_{X}\varepsilon \varepsilon ^{\prime })] \\
=\frac{1}{n}tr(\mathbb{E[}M_{X}\mathbb{E}[\varepsilon \varepsilon ^{\prime
}|X]])=\frac{\sigma ^{2}}{n}\mathbb{E[}tr(M_{X})] \\
=\frac{\sigma ^{2}}{n}\mathbb{E[}tr(I_{n}-P_{X})]=\frac{\sigma ^{2}}{n}%
\mathbb{E}\left[ tr(I_{n})-tr(X(X^{\prime }X)^{-1}X^{\prime })\right] \\
=\frac{\sigma ^{2}}{n}\mathbb{E}\left[ n-tr(X^{\prime }X(X^{\prime }X)^{-1})%
\right] =\frac{\sigma ^{2}}{n}\mathbb{E}\left[ n-tr(I_{k})\right] =\frac{%
\sigma ^{2}(n-k)}{n}\text{.}
\end{gather*}%
The bias is given by 
\begin{equation*}
\left( \mathbb{E}[\widehat{\sigma }_{n}^{2}]-\sigma ^{2}\right) =\frac{%
\sigma ^{2}(n-k)}{n}-\sigma ^{2}=\frac{-k\sigma ^{2}}{n}
\end{equation*}%
so estimator is negatively biased. As $n\rightarrow \infty ,$ $\frac{%
-k\sigma ^{2}}{n}\rightarrow 0.$

\bigskip

Question 3

Consider the linear regression model from Question 2. Recall that the ridge
estimator is defined as $\hat{\beta}^{Ridge}=\left( X^{\prime }X+\lambda
_{n}I_{k}\right) ^{-1}X^{\prime }y.$ Suppose that $\lambda _{n}=cn^{\delta }$
for some $\delta \geq 0.$ For which range of values for $c$ and $\delta $ is 
$\hat{\beta}^{Ridge}$ a consistent estimator for $\beta ?$ [15 marks] 
\begin{eqnarray*}
p\lim_{n\rightarrow \infty }\hat{\beta}^{Ridge} &=&\left(
p\lim_{n\rightarrow \infty }\frac{1}{n}\tsum\nolimits_{i=1}^{n}\mathbf{x}_{i}%
\mathbf{x}_{i}^{\prime }+\lim_{n\rightarrow \infty }cn^{\delta
-1}I_{k}\right) ^{-1}p\lim_{n\rightarrow \infty }\frac{1}{n}%
\tsum\nolimits_{i=1}^{n}\mathbf{x}_{i}\mathbf{x}_{i}^{\prime }\beta  \\
&&+\left( p\lim_{n\rightarrow \infty }\frac{1}{n}\tsum\nolimits_{i=1}^{n}%
\mathbf{x}_{i}\mathbf{x}_{i}^{\prime }+\lim_{n\rightarrow \infty }cn^{\delta
-1}I_{k}\right) ^{-1}p\lim_{n\rightarrow \infty }\frac{1}{n}%
\tsum\nolimits_{i=1}^{n}\mathbf{x}_{i}\mathbf{\varepsilon }_{i}
\end{eqnarray*}

so for any values of $c\in \left( -\infty ,\infty \right) $ and any $\delta
\in \lbrack 0,1),$ $\lim_{n\rightarrow \infty }cn^{\delta -1}I_{k}=0$ so $%
p\lim_{n\rightarrow \infty }\hat{\beta}^{Ridge}=\beta .$

\end{document}
